% !TEX root = ../main.tex
\section{TSGD as a Building Block of Modern Optimizers}
\label{sec:connections}

\subsection{Transport Is Already Everywhere}
\label{subsec:refreshing-and-transporting}

Many modern stochastic optimizers combine fresh gradients with
transported estimates. Their common building blocks are the three
operations in Figure~\ref{fig:fundamental-estimators}: fresh estimation,
averaging, and transport. Write $H_{k+1}$ for a fresh unbiased estimate of
$\nabla f(x_{k+1})$, $D_{k+1}$ for an estimate of its change from $x_k$,
and $g_k+D_{k+1}$ for the transported estimate. A stochastic gradient
difference uses the same fresh sample or minibatch at both endpoints.
Table~\ref{tab:transport-building-blocks} summarizes the resulting
estimator relationships. Here $\alpha,\beta\in[0,1]$ weight fresh
estimates, and $\gamma$ scales the MARS transport correction.

\begin{table}[htbp]
\centering
\small
\setlength{\tabcolsep}{4pt}
\renewcommand{\arraystretch}{1.3}
\begin{tabular}{@{}>{\raggedright\arraybackslash}p{0.20\textwidth}
                  >{\raggedright\arraybackslash}p{0.43\textwidth}
                  >{\raggedright\arraybackslash}p{0.31\textwidth}@{}}
\toprule
Method & Estimator at iteration $k+1$ & Role of transport \\
\midrule
SARAH / SPIDER
& $H_{k+1}$ or $g_k+D_{k+1}$
& Transport between periodic refreshes. \\
\addlinespace[3pt]
Geom-SARAH
& $H_{k+1}$ or $g_k+D_{k+1}$
& Transport for a geometric number of steps. \\
\addlinespace[3pt]
PAGE
& $H_{k+1}$ or $g_k+D_{k+1}$
& Refresh with probability $p$. \\
\addlinespace[3pt]
MARINA
& $\nabla f(x_{k+1})$ or $g_k+D_{k+1}^{\mathcal Q}$
& Communicate compressed gradient changes. \\
\addlinespace[3pt]
STORM / MVR
& $\begin{aligned}[t]
   &(1-\alpha)(g_k+D_{k+1})\\
   &\quad+\alpha H_{k+1}
   \end{aligned}$
& Average fresh and transported estimates. \\
\addlinespace[3pt]
MARS (unclipped)
& $\begin{aligned}[t]
   &(1-\beta)(g_k+\gamma D_{k+1})\\
   &\quad+\beta H_{k+1}
   \end{aligned}$
& Scale the transport correction within momentum. \\
\bottomrule
\end{tabular}
\caption{Transport in modern optimizers. Rows with two alternatives
select one branch according to a restart or refresh rule; the last two
rows form weighted combinations. $D_{k+1}^{\mathcal Q}$ denotes an
unbiased compressed increment. The identities describe estimators;
batch sizes, stepsizes, and iterate updates belong to each algorithm.}
\label{tab:transport-building-blocks}
\end{table}

SVRG \citep{NIPS2013_ac1dd209} and SAGA
\citep{defazio2014sagafastincrementalgradient} also use stochastic
gradient differences to reduce variance. SVRG compares component
gradients at the current iterate and a reference point held fixed
within each epoch, adding the full gradient at that reference point.
SAGA instead stores past gradients for individual components and
combines the sampled component's gradient difference with the average
of the stored gradients, updating one entry at each iteration.
Both exploit noise cancellation in gradient differences. Unlike TSGD,
they do not transport the previous gradient estimate along successive
iterates: each estimate combines a current component gradient with
stored reference information. In their standard forms, these estimates
are conditionally unbiased, whereas pure TSGD inherits and accumulates
tracking errors.

SARAH and SPIDER transport an estimate between full-gradient or
large-batch recomputations
\citep{nguyen2017sarah,nguyen2017sarah_nonconvex,fang2018spider}.
Inexact SARAH allows batch initialization \citep{nguyen2021inexact},
while Geom-SARAH randomizes the phase length
\citep{horvath2022adaptivity}. PAGE refreshes probabilistically
\citep{li2021page}; MARINA combines exact refreshes with compressed
gradient differences \citep{gorbunov2021marina}.

MVR/STORM instead averages fresh and transported estimates at each
iteration \citep{cutkosky2019momentum}. The unclipped MARS estimator
\citep{yuan2025mars} includes MSGD at $\gamma=0$ and MVR at $\gamma=1$.
These identities concern the stated sampling: cached gradients,
clipping, and preconditioning require separate treatment. Sampling
also affects the noise covariance, as emphasized by the VENUS
construction \citep{venus2026draft}.
Appendix~\ref{app:algorithmic-details} gives the precise identities,
error recursions, and special cases.

Gradient differences also appear in decentralized gradient tracking
\citep{Qu_2018}, where agents combine changes in their local gradients
with information from neighboring agents to track the average gradient
across the network.

\subsection{Why Combine Fundamental Estimators?}
\label{subsec:why-combine-estimators}

SGD, MSGD, and TSGD use fresh information, averaging, and transport,
respectively. Modern optimizers combine them through weighted averages
or alternate between fresh estimation and transport. Three motivations
explain why these combinations can be useful.

\emph{Computational efficiency.} An accurate absolute gradient may
require a large batch, whereas a small batch can suffice for an accurate
gradient change. Even with two evaluations per sample, transport can
then cost much less. SARAH, SPIDER, and PAGE preserve information from
an expensive gradient estimate through many cheaper transport steps.
Refreshing more often limits accumulated error but reduces this saving.

\emph{Complementary strengths.} Momentum averages noisy information
from previous iterations, but those gradients were evaluated at points
the iterate has already left. Transport accounts for the change in
the true gradient, reducing tracking lag while retaining the benefit
of averaging. Fresh gradients supply absolute information that pure
transport cannot recover: mixing them in damps initialization and
accumulated transport errors, while introducing new pointwise noise.
MVR, STORM, and MARS balance these roles through their mixing weights
and transport corrections. This complementarity matters even when
the noise regime remains fixed.

\emph{Changing regimes.} The relative accuracy of fresh and transported
estimates can change during a single run. Similarity bounds each new
transport error in proportion to the squared displacement, so shorter
steps tighten this bound without removing previously accumulated error.
Section~\ref{subsec:fixed-variance-variable-similarity} illustrates
this displacement dependence even when pointwise noise variance stays
constant. Fresh gradients may also become more reliable: in
Section~\ref{subsec:multiplicative-noise-comparison}, their noise
vanishes near the minimizer, while pure transport can retain a nonzero
tracking error. Thus the balance can change as iterates move, steps
shrink, and errors accumulate.

Mixing weights and refresh frequency determine how much an optimizer
relies on each mechanism. Fixed choices provide a compromise, while
adapting them can respond to changing conditions. The balance matters:
combining estimators does not automatically improve convergence, and
a poor choice can perform worse than using one estimator alone.
